\documentclass[11pt,a4paper]{article}
\usepackage[margin=1in]{geometry}
\usepackage{amsmath,amssymb,amsthm}
\usepackage{algorithm}
\usepackage{algpseudocode}
\usepackage{booktabs}
\usepackage{hyperref}

\theoremstyle{definition}
\newtheorem{definition}{\textbf{Definition}}[section]
\newtheorem{theorem}{\textbf{Theorem}}[section]
\newtheorem{lemma}[theorem]{\textbf{Lemma}}
\newtheorem{remark}{\textbf{Remark}}[section]

\title{\textbf{Expert Parallelism in Mixture-of-Experts: All-to-All Dispatch and Combine Routing}}
\author{\textbf{Team Scaibu} \\ \small Email: \texttt{jaydeep.9664920749@gmail.com} \\ \small \textbf{Architecture and Core Engine Team}}
\date{\textbf{October 2026}}

\begin{document}
\maketitle

\section{Definitions and Cognitive Mental Models}

\subsection{Formal Mathematical Definition}
\textbf{Definition (Expert Parallelism and All-to-All Dispatch):}
In sparse Mixture-of-Experts (MoE) architectures, the total parameter count scales to hundreds of billions or trillions of parameters by replacing standard feed-forward networks (FFNs) with an ensemble of $E$ specialized expert networks $\{F_1, \dots, F_E\}$. Because no single GPU can host all $E$ experts, \textbf{Expert Parallelism (EP)} distributes the $E$ experts across $N_{\text{ep}}$ distinct device ranks:
{\boldmath
\begin{equation}
\mathcal{E}_r = \left\{ F_e \mid \operatorname{rank}(e) = r \right\}, \quad |\mathcal{E}_r| = \frac{E}{N_{\text{ep}}}
\end{equation}
}
During execution, a routing gating network $G(x) = \operatorname{TopK}(\operatorname{Softmax}(W_g x), k)$ selects the top-$k$ experts for each token. Because tokens on device $i$ frequently choose experts hosted on remote device $j$ ($j \neq i$), execution requires a collective \textbf{All-to-All Dispatch} exchange:
{\boldmath
\begin{equation}
X_{\text{local}} \xrightarrow{\text{All-to-All Dispatch}} X_{\text{routed}}, \quad X_{\text{routed}} \in \mathbb{R}^{C_{\text{expert}} \times d}
\end{equation}
}
Each device executes its locally hosted experts on the tokens routed to it. A second \textbf{All-to-All Combine} collective returns the computed expert output activations back to the tokens' originating devices, where they are weighted by gating coefficients:
{\boldmath
\begin{equation}
y = \sum_{e \in \operatorname{TopK}} g_e F_e(x)
\end{equation}
}
The total network data transferred across both All-to-All phases per MoE layer is:
{\boldmath
\begin{equation}
\text{Volume}_{\text{All-to-All}} = 2 \cdot (B \cdot L \cdot k) \cdot d \cdot B_{\text{bytes}}
\end{equation}
}
which acts as the primary throughput bottleneck in distributed MoE systems.

\subsection{Expanded Non-Technical Definition (Intuition for Anyone)}
\begin{enumerate}
    \item \textbf{A Global Hospital with 64 Specialty Clinics}: Each GPU is a specialty clinic: GPU 1 has 8 Brain Surgeons, GPU 2 has 8 Heart Surgeons, etc. Incoming patients (tokens) arrive at any local clinic.
    \item \textbf{The Dispatch Airport (All-to-All)}: The local clinic triages each patient and puts them on high-speed jets to their chosen specialty clinics. GPU 1 sends heart patients to GPU 2 and receives brain patients from GPU 2 simultaneously.
    \item \textbf{The Return Flight (Combine)}: Once the specialists finish their surgery, the patients fly back to their original hospital with their treatment results in hand.
    \item \textbf{The Straggler Trap (Load Imbalance)}: If 90\% of patients suddenly want to see the Heart Surgeon, GPU 2 gets overwhelmed and has a massive waiting line, while GPU 1 sits completely empty! Load balancing algorithms ensure patients are distributed evenly across specialists.
\end{enumerate}

\subsection{Child-Friendly Mental Model (Physical Metaphor)}
Imagine a massive school project:
\begin{itemize}
    \item \textbf{Table 1} has all the colored markers.
    \item \textbf{Table 2} has all the scissors.
    \item \textbf{Table 3} has all the glue sticks.
    \item All the kids trade places at the exact same second: the kids who need scissors run to Table 2, the kids who need glue run to Table 3, and then they all trade back!
\end{itemize}

\section{Core Mathematical Equations and Definitions}

\subsection{All-to-All Communication Volume}
{\boldmath
\begin{equation}
V_{\text{MoE}} = 2 \cdot B \cdot L \cdot k \cdot d \cdot B_{\text{bytes}}
\end{equation}
}
\begin{itemize}
    \item \textbf{Formal Definition}: Total bidirectional bytes transferred per MoE layer across all cluster ranks.
    \item \textbf{What It Means}: Two passes (dispatch and combine) carrying $k$ activation vectors per sequence token.
    \item \textbf{Why It Is Important}: Primary metric governing MoE network bandwidth requirements.
\end{itemize}

\subsection{Expert Load Imbalance Ratio}
{\boldmath
\begin{equation}
\mathcal{I}_{\text{load}} = \frac{\max_r \sum_{e \in \mathcal{E}_r} C_e}{\frac{1}{N_{\text{ep}}} \sum_{r=1}^{N_{\text{ep}}} \sum_{e \in \mathcal{E}_r} C_e}
\end{equation}
}
\begin{itemize}
    \item \textbf{Formal Definition}: Ratio of maximum rank token workload to average rank token workload.
    \item \textbf{What It Means}: $\mathcal{I}_{\text{load}} = 1.0$ indicates perfect load balance; values $>1.5$ cause severe hardware stalling.
    \item \textbf{Why It Is Important}: Core metric monitored by auxiliary load-balancing loss terms.
\end{itemize}

\section{Rigorous Mathematical Analysis Checklist}

\begin{itemize}
    \item \textbf{Notation}: $E$ (total experts), $N_{\text{ep}}$ (expert ranks), $k$ (top-k active experts), $d$ (hidden dimension), $C_e$ (token count per expert).
    \item \textbf{Epistemic Status}: Fedus et al. (Switch Transformer, 2022); Lepikhin et al. (GShard, 2021). Universal MoE serving standard.
    \item \textbf{Dependency Graph}: Token $x \to \text{Gating Router } \to \text{All-to-All Dispatch} \to \text{Expert FFN} \to \text{All-to-All Combine} \to y$.
    \item \textbf{Dimensions}: Total experts $E \in [8, 512]$, active experts $k \in [1, 8]$, hidden dimension $d \in [2048, 8192]$.
    \item \textbf{Regimes}: Frontier sparse models (e.g. Mixtral 8x7B, DBRX, Grok, DeepSeek-V2/V3).
    \item \textbf{Invariants}: Total tokens dispatched across all ranks equals exactly $B \cdot L \cdot k$.
    \item \textbf{Isomorphisms}: Isomorphic to complete bipartite packet routing and distributed hash table (DHT) key shuffling.
\end{itemize}

\section{Theoretical Theorems and Formal Proofs}

\begin{theorem}[Conservation of Dispatched Token Mass]
Let a global batch consist of $B$ sequences of length $L$, and let each token select exactly $k$ distinct experts. For any arbitrary expert-to-rank assignment mapping $\phi: \{1, \dots, E\} \to \{1, \dots, N_{\text{ep}}\}$, the sum of tokens received across all $N_{\text{ep}}$ destination ranks is invariant and strictly equal to $B \cdot L \cdot k$:
{\boldmath
\begin{equation}
\sum_{r=1}^{N_{\text{ep}}} \operatorname{TokensReceived}(r) \equiv B \cdot L \cdot k
\end{equation}
}
\end{theorem}

\begin{proof}
Let $T = B \cdot L$ denote the total number of input tokens across all ranks.
For each token $t \in \{1, \dots, T\}$, the gating router selects a set of experts $S_t \subset \{1, \dots, E\}$ with $|S_t| = k$.
The total number of expert activations instantiated across the entire cluster is:
{\boldmath
\begin{equation}
N_{\text{total\_activations}} = \sum_{t=1}^T |S_t| = \sum_{t=1}^T k = T \cdot k = B \cdot L \cdot k
\end{equation}
}
Each expert activation $(t, e)$ for $e \in S_t$ is routed to destination rank $r = \phi(e)$.
Let indicator $\mathbb{I}(\phi(e) = r) = 1$ if expert $e$ resides on rank $r$, and $0$ otherwise.
The tokens received by rank $r$ is:
{\boldmath
\begin{equation}
\operatorname{TokensReceived}(r) = \sum_{t=1}^T \sum_{e \in S_t} \mathbb{I}(\phi(e) = r)
\end{equation}
}
Summing over all ranks $r \in \{1, \dots, N_{\text{ep}}\}$:
{\boldmath
\begin{align}
\sum_{r=1}^{N_{\text{ep}}} \operatorname{TokensReceived}(r) &= \sum_{r=1}^{N_{\text{ep}}} \sum_{t=1}^T \sum_{e \in S_t} \mathbb{I}(\phi(e) = r) = \sum_{t=1}^T \sum_{e \in S_t} \left( \sum_{r=1}^{N_{\text{ep}}} \mathbb{I}(\phi(e) = r) \right)
\end{align}
}
Since $\phi$ is a well-defined function, every expert $e$ resides on exactly one rank, so $\sum_{r=1}^{N_{\text{ep}}} \mathbb{I}(\phi(e) = r) = 1$.
Therefore:
{\boldmath
\begin{equation}
\sum_{r=1}^{N_{\text{ep}}} \operatorname{TokensReceived}(r) = \sum_{t=1}^T \sum_{e \in S_t} 1 = \sum_{t=1}^T k = B \cdot L \cdot k
\end{equation}
}
This establishes exact mass conservation across all possible routing trajectories.
\end{proof}

\begin{theorem}[Lower Bound on Cluster Execution Time Under Load Imbalance]
Let $t_{\text{expert}}$ be the compute duration per token on an expert. In a cluster of $N_{\text{ep}}$ expert-parallel ranks, the compute makespan of an MoE layer is strictly lower-bounded by the maximum rank load:
{\boldmath
\begin{equation}
T_{\text{layer}} \ge t_{\text{expert}} \cdot \max_{r} \operatorname{TokensReceived}(r) = t_{\text{expert}} \cdot \mathcal{I}_{\text{load}} \cdot \left( \frac{B \cdot L \cdot k}{N_{\text{ep}}} \right)
\end{equation}
}
\end{theorem}

\begin{proof}
Because synchronization barriers occur at the conclusion of each MoE layer before subsequent attention layers can begin, the total layer duration is governed by the slowest worker rank (the maximum load rank):
{\boldmath
\begin{equation}
T_{\text{layer}} = \max_{r \in [1, N_{\text{ep}}]} T_{\text{rank}}(r) = t_{\text{expert}} \cdot \max_r \operatorname{TokensReceived}(r)
\end{equation}
}
By definition of the load imbalance ratio $\mathcal{I}_{\text{load}}$:
{\boldmath
\begin{equation}
\max_r \operatorname{TokensReceived}(r) = \mathcal{I}_{\text{load}} \cdot \operatorname{MeanTokens} = \mathcal{I}_{\text{load}} \cdot \left( \frac{\sum_{r=1}^{N_{\text{ep}}} \operatorname{TokensReceived}(r)}{N_{\text{ep}}} \right)
\end{equation}
}
Substituting the conservation theorem result $\sum_r \operatorname{TokensReceived}(r) = B \cdot L \cdot k$:
{\boldmath
\begin{equation}
T_{\text{layer}} \ge t_{\text{expert}} \cdot \mathcal{I}_{\text{load}} \cdot \left( \frac{B \cdot L \cdot k}{N_{\text{ep}}} \right)
\end{equation}
}
If $\mathcal{I}_{\text{load}} = 2.0$ (one rank receives twice its fair share), cluster execution time doubles, wasting 50\% of total cluster compute capacity.
\end{proof}

\section{Foundational Architectural Questions and Epistemic Meta-Questions}

\subsection{Architectural Question 1: Expert Capacity Limits and Token Dropping}
\textbf{Question:} How do systems prevent GPUs from running out of memory when too many tokens choose the same expert?\\
\textbf{Answer:} Systems enforce an \textbf{Expert Capacity Factor} $C_{\text{factor}} \in [1.0, 1.5]$, setting maximum tokens per expert to $C = C_{\text{factor}} \cdot \frac{B \cdot L \cdot k}{E}$. If more than $C$ tokens route to expert $e$, excess tokens are dropped (bypassed around the expert via residual connection without modification).

\subsection{Architectural Question 2: Hierarchical All-to-All Routing}
\textbf{Question:} How is All-to-All communication optimized across thousands of GPUs?\\
\textbf{Answer:} Instead of running a flat All-to-All across all nodes, modern architectures (DeepSeek-V3, Mixtral) deploy \textbf{Hierarchical MoE}:
1. Intra-node All-to-All over high-speed NVLink.
2. Inter-node aggregation over InfiniBand networks.
Tokens are pre-aggregated locally before crossing inter-node switches, reducing expensive network hops by up to $4\times$.

\section{Algorithmic Implementation Specification}

\begin{algorithm}[H]
\caption{Expert Parallelism All-to-All Dispatch Accounting}
\begin{algorithmic}[1]
\Require Token expert selections $\mathbf{T} \in \mathbb{N}^{B \times k}$, expert-to-rank map $\phi$, rank count $N$, dim $d$, byte size $B_{\text{byte}}$
\Ensure Rank counts $\mathbf{C} \in \mathbb{N}^N$, comm bytes $V$, imbalance $\mathcal{I}$, total routed $T_{\text{routed}}$
\State Allocate $\mathbf{C} \gets [0]^N, \quad T_{\text{routed}} \gets 0$
\For{$t = 0$ \textbf{to} $B-1$}
    \For{$idx \in \mathbf{T}[t]$}
        \State $r \gets \phi(idx) \pmod N$
        \State $\mathbf{C}[r] \gets \mathbf{C}[r] + 1$
        \State $T_{\text{routed}} \gets T_{\text{routed}} + 1$
    \EndFor
\EndFor
\State $C_{\max} \gets \max_r \mathbf{C}[r]$
\State $C_{\text{mean}} \gets T_{\text{routed}} / N$
\State $\mathcal{I} \gets C_{\max} / \max(10^{-12}, C_{\text{mean}})$
\State $V \gets 2 \cdot T_{\text{routed}} \cdot d \cdot B_{\text{byte}}$ \Comment{Dispatch + Combine passes}
\State \Return $(\mathbf{C}, V, \mathcal{I}, T_{\text{routed}})$
\end{algorithmic}
\end{algorithm}

\section{Big-$\Theta$ Complexity Analysis}

\begin{itemize}
    \item \textbf{Routing Classification Time}: $\Theta(B \cdot k)$ hash and count operations.
    \item \textbf{All-to-All Network Transfer Volume}: $\Theta(B \cdot k \cdot d)$ total elements.
    \item \textbf{Memory Storage per Device}: $\Theta\left(\frac{E}{N_{\text{ep}}} \cdot P_{\text{expert}}\right)$ local weights.
    \item \textbf{Parallel Speedup}: Linear $\Theta(N_{\text{ep}})$ reduction in per-node expert memory.
\end{itemize}

\section{Single Canonical Repository Reference}

The complete deterministic scalar implementation, verification suite, and unit test benchmarks are published at:
\begin{center}
\url{https://github.com/Chief-Strategist-J/llm-obs-infra}
\end{center}

\end{document}
